% !TeX root = ../教材题库_全部条目_转写工作稿.tex
% Source: XJ-B1-C04-TH-P115-Y0591
% 第4章 / 4.3 对数函数 / 教材 p.115 / PDF 第122页
% OCR未校对：公式与文字均可能有误；未完成数学TeX转写；禁止直接发给学生
% 原题及相邻原解截图（核对用）：../题目原图/XJ-B1-C04-TH-P115-Y0591-01.webp
% 本文件只在转写完成后作为题干使用；OCR 条目尚不是可靠题干。
\begin{bankitem}{p.115 思考（含紧邻说明）}{第4章 / 4.3 对数函数 / 教材 p.115 / PDF 第122页}{XJ-B1-C04-TH-P115-Y0591}{OCR未校对：公式与文字均可能有误；未完成数学TeX转写；禁止直接发给学生}
\begin{ocrdraft}
[教材p.115 / PDF第122页]
\par
？思考我们在4．2．1小节的例2中知道，钴60射线穿过厚度为1cm的铅板
\par
后，强度变为原来的0．568倍，穿过厚度为工cm的铅板后的强度与原来的强度之比
\par
为H(x）一0．568能若要使射线穿过铅板后的强度是原来的百分之一，铅板厚度应为
\par
多少呢？也就是说，知道568J=0．01，如何求工呢？
\par
把上述问题抽象化为一般的数学问题：在等式
\par
N=>0且a1）
\par
中，知道了N和。需求出b的值．
\par
由8=2容易看出乃一3，但由2一10b求b就不容易了．
\par
这是一种新的数学运算，叫作对数运算．具体来说：
\par
对数的发明者是苏格
\par
如果一N（>0且1，N>0），那么叫作以为底，
\par
兰数学家纳皮尔．“嵊”是
\par
正数N的对数，记作b=logN.这里，磊叫作对数的底数，N
\par
拉丁文logarithmus(对数）
\par
叫作对数的真数．
\par
的缩写．
\end{ocrdraft}
\end{bankitem}
